\begin{frame}[plain]\FrameAnchor{V23-title}
\centering
{\Large\bfseries\color{courseblue}第 23 卷\par}
\vspace{0.35cm}
{\LARGE\bfseries 粒子物理、标准模型与强子量子数\par}
\vspace{0.35cm}
{\large 补齐从规范群、粒子谱、离散对称性到散射与因子化的粒子物理背景。\par}
\vfill
\CourseID{Tbl}{23.00}\quad 5 个学习单元，固定编号不随合订方式改变
\end{frame}

\begin{frame}[t]{第 23 卷路线图}\FrameAnchor{V23-route}
\small
\begin{tabularx}{\linewidth}{p{0.14\linewidth}Y}
\toprule 编号 & 学习单元\\\midrule
23.01 & 一代标准模型场表、手征性与电荷\\
23.02 & 夸克模型、flavor 多重态与强子分类\\
23.03 & 角动量耦合、宇称、荷共轭与 JPC\\
23.04 & 两体相空间、散射截面与衰变宽度\\
23.05 & 显式卷积、硬散射因子化与部分子模型\\
\bottomrule\end{tabularx}
\vfill
\SourceLine{BOOK-QFT；BOOK-GROUP；BOOK-LQCD}
\end{frame}

\begin{frame}[t]{先修诊断与本卷出口}\FrameAnchor{V23-diagnostic}
\begin{columns}[T,onlytextwidth]
\column{0.48\textwidth}\begin{block}{开始前应能回答}
\small\begin{itemize}
\item V07
\item V16
\item V22
\end{itemize}\end{block}
\column{0.48\textwidth}\begin{block}{学完后应能独立完成}
\small\begin{itemize}
\item 能读标准模型量子数表
\item 能用 JPC 和 flavor 分类强子
\item 能连接散射振幅、截面和部分子因子化
\end{itemize}\end{block}
\end{columns}
\vfill\scriptsize 若先修项答不出，回到相应前卷；不要以术语熟悉代替可计算能力。
\end{frame}

\begin{frame}[t]{定量图：两体相空间阈值}\FrameAnchor{V23-chart}
\centering\begin{tikzpicture}
\begin{axis}[width=0.88\linewidth,height=0.63\textheight,xmin=0,xmax=4,ymin=0,ymax=2.1,xlabel={$m_1+m_2$ 以上的超额能量（归一化）},ylabel={$p^\ast$（示意）},grid=major,grid style={courseline!55},legend style={font=\scriptsize,draw=none,fill=white},tick label style={font=\scriptsize},label style={font=\small}]
\addplot[very thick,courseblue,domain=0:4,samples=160] {sqrt(x)};
\addlegendentry{S 波相空间}
\end{axis}\end{tikzpicture}
\par\scriptsize\FigureID{23.00}：非相对论阈值 p*\ensuremath{\propto}\ensuremath{\sqrt{E}}；实际系数由质量决定。
\SourceLine{两体运动学}
\end{frame}

\begin{frame}[t]{一代标准模型场表、手征性与电荷：问题与图像}\FrameAnchor{23.01-1}
\footnotesize
\begin{block}{本节问题}怎样从 (SU(3)c,SU(2)L)Y 表示标签逐个读出一代粒子的电荷？\end{block}
\PhysicalFlow{先标左右手 Weyl 场}{记录三个规范表示}{用 Q=T3+Y\_EW 核对每个分量}{23.01}
\begin{block}{物理原则}\KnowledgeID{23.01}\quad 做反常计算时把右手场改写成左手共轭 uR\ensuremath{^{c}}、dR\ensuremath{^{c}}、eR\ensuremath{^{c}}：非阿贝尔表示取共轭且 Y\_EW 反号；读取任何文献表前先确认其是否使用 Q=T3+Y/\allowbreak{}2。\end{block}
\SourceLine{BOOK-QFT}
\end{frame}

\begin{frame}[t]{一代标准模型场表、手征性与电荷：定义与主公式}\FrameAnchor{23.01-2}
\footnotesize\begin{tabularx}{\linewidth}{p{0.30\linewidth}Y}
\toprule 术语 & 本课程中的精确定义\\\midrule
\DefinitionID{23.01-f178d52e23} 手征性 & Lorentz 旋量的左/\allowbreak{}右投影；SU(2)L 只作用在左手 doublet\\
\DefinitionID{23.01-33530d16a6} 表示标签 & (Rc,Rw)Y\_EW 分别记录颜色表示、弱同位旋表示和电弱超荷\\
\DefinitionID{23.01-286ecdb958} 电弱超荷 Y\_EW & 本课采用 Q=T3+Y\_EW；不同于强子分类的 Y\_flavor\\
\bottomrule\end{tabularx}\vspace{0.15cm}
\begin{block}{\EquationID{23.01} 主公式}\centering\CourseDisplayEquation{Q_L:(\mathbf3,\mathbf2)_{1/6},\ u_R:(\mathbf3,\mathbf1)_{2/3},\ d_R:(\mathbf3,\mathbf1)_{-1/3},\ L_L:(\mathbf1,\mathbf2)_{-1/2},\ e_R:(\mathbf1,\mathbf1)_{-1};\quad Q=T_3+Y_{\rm EW}}\end{block}
这五个 multiplet 加 Higgs H:(1,2)1/\allowbreak{}2 构成无右手中微子的一代最小场表；doublet 上下分量取 T3=\ensuremath{\pm}1/\allowbreak{}2，singlet 取 T3=0。
\end{frame}

\begin{frame}[t]{一代标准模型场表、手征性与电荷：推导与证据边界}\FrameAnchor{23.01-3}
\footnotesize
\DerivationID{23.01}\quad 由本节定义与前置结论逐步推出 \CourseRef{Eq}{23.01}：
\begin{enumerate}
\item Q\_L 的上分量电荷 1/\allowbreak{}2+1/\allowbreak{}6=2/\allowbreak{}3，下分量 -1/\allowbreak{}2+1/\allowbreak{}6=-1/\allowbreak{}3。
\item L\_L 给 Q\ensuremath{\nu}=1/\allowbreak{}2-1/\allowbreak{}2=0、Qe=-1/\allowbreak{}2-1/\allowbreak{}2=-1；右手 singlet 的电荷直接等于 Y\_EW。
\item H 的上、下分量电荷为 +1、0；中性下分量可取真空期望值。Yukawa 单项的颜色、SU(2) 缩并和 Y\_EW 总和都必须为 singlet/\allowbreak{}零。
\end{enumerate}\begin{block}{可执行检查}
\SympyID{23.01}\quad 一代标准模型场表、手征性与电荷；2/2 项通过；SymPy exact。
\par\scriptsize 假设：采用 Q=T3+Y、左手夸克双重态 Y=1/\allowbreak{}6
\end{block}
\BoundaryLine{只核对电荷归一化；规范反常与完整粒子谱另验。}
\end{frame}

\begin{frame}[t]{一代标准模型场表、手征性与电荷：计算算法}\FrameAnchor{23.01-4}
\footnotesize
\AlgorithmID{23.01}\begin{enumerate}
\item 为每个场保存 chirality、Rc、Rw、Y\_EW 四列，禁止把颜色、flavor 与弱指标复用。
\item 自动枚举 doublet 的 T3 并核对电荷。
\item 转成全左手表时共轭表示并反转所有 Abelian 荷。
\item 逐项检查 Qbar\_L H d\_R、Qbar\_L Htilde u\_R、Lbar\_L H e\_R 的表示和超荷守恒。
\end{enumerate}\begin{columns}[T,onlytextwidth]
\column{0.56\textwidth}\begin{block}{停止条件与交叉检查}\scriptsize\begin{itemize}
\item[\CheckMark] H 中性分量满足 T3+Y\_EW=-1/\allowbreak{}2+1/\allowbreak{}2=0。
\item[\CheckMark] 粒子改为左手反粒子时电荷和 Y\_EW 反号。
\end{itemize}\end{block}
\column{0.40\textwidth}\begin{block}{代码映射}
\scriptsize \texttt{课程内纸笔/\allowbreak{}符号计算练习}
\end{block}\end{columns}\end{frame}

\begin{frame}[t]{一代标准模型场表、手征性与电荷：练习与完整解答}\FrameAnchor{23.01-5}
\footnotesize
\begin{block}{\ExerciseID{23.01}}把 u\_R 改写成左手共轭场 u\_R\ensuremath{^{c}}；它的 (SU(3)c,SU(2)L)Y\_EW 标签和电荷是什么？\end{block}
\begin{exampleblock}{\SolutionID{23.01}}u\_R 原为 (3,1)2/\allowbreak{}3；共轭后为 (3bar,1)-2/\allowbreak{}3，T3=0，所以左手反粒子场的电荷为 -2/\allowbreak{}3。\end{exampleblock}
\vfill
\scriptsize 回看：\CourseRef{Def}{23.01-f178d52e23}；\CourseRef{Eq}{23.01}；\CourseRef{SYM}{23.01}；\CourseRef{Alg}{23.01}。
\SourceLine{BOOK-QFT}
\end{frame}

\begin{frame}[t]{夸克模型、flavor 多重态与强子分类：问题与图像}\FrameAnchor{23.02-1}
\footnotesize
\begin{block}{本节问题}质量谱中的八重态、十重态和介子 nonet 从何而来？\end{block}
\PhysicalFlow{轻 flavor SU(3)}{张量积 3\ensuremath{\otimes}3bar 或 3\ensuremath{^{3}}}{按 I、Y\_flavor、S 排列多重态}{23.02}
\begin{block}{物理原则}\KnowledgeID{23.02}\quad flavor SU(3) 被 ms\ensuremath{\ne}mu,d 显著破坏；多重态用于分类和近似关系，不代表精确简并。\end{block}
\SourceLine{BOOK-QFT；BOOK-GROUP；BOOK-LQCD}
\end{frame}

\begin{frame}[t]{夸克模型、flavor 多重态与强子分类：定义与主公式}\FrameAnchor{23.02-2}
\footnotesize\begin{tabularx}{\linewidth}{p{0.30\linewidth}Y}
\toprule 术语 & 本课程中的精确定义\\\midrule
\DefinitionID{23.02-8fa457f743} 同位旋 & 近似 u,d flavor SU(2) 对称的量子数\\
\DefinitionID{23.02-1493fa31b2} 奇异数 S & 强子中反 s 数减 s 数；s 夸克贡献 S=-1\\
\DefinitionID{23.02-0fb80e2e21} 强子超荷 Y\_flavor & Y\_flavor=B+S，用于 flavor 权图，绝不能与电弱 Y\_EW 混用\\
\bottomrule\end{tabularx}\vspace{0.15cm}
\begin{block}{\EquationID{23.02} 主公式}\centering\CourseDisplayEquation{\mathbf3\otimes\mathbf3\otimes\mathbf3=\mathbf{10}\oplus\mathbf8\oplus\mathbf8\oplus\mathbf1}\end{block}
三个 flavor triplet 的 27 维空间分成全对称十重态、两个混合对称八重态和全反对称 singlet。
\end{frame}

\begin{frame}[t]{夸克模型、flavor 多重态与强子分类：推导与证据边界}\FrameAnchor{23.02-3}
\footnotesize
\DerivationID{23.02}\quad 由本节定义与前置结论逐步推出 \CourseRef{Eq}{23.02}：
\begin{enumerate}
\item 先分解 3\ensuremath{\otimes}3=6\_s\ensuremath{\oplus}3bar\_a。
\item 再乘 3：6\ensuremath{\otimes}3=10\ensuremath{\oplus}8，3bar\ensuremath{\otimes}3=8\ensuremath{\oplus}1。
\item 维数核对 10+8+8+1=27。
\end{enumerate}\begin{block}{可执行检查}
\SympyID{23.02}\quad 夸克模型、flavor 多重态与强子分类；1/1 项通过；SymPy exact。
\par\scriptsize 假设：三夸克 flavor SU(3) 张量积
\end{block}
\BoundaryLine{维数相等是必要检查，不单独证明不可约分解。}
\end{frame}

\begin{frame}[t]{夸克模型、flavor 多重态与强子分类：计算算法}\FrameAnchor{23.02-4}
\footnotesize
\AlgorithmID{23.02}\begin{enumerate}
\item 由 flavor 组成生成候选强子算符。
\item 结合颜色全反对称和费米统计确定自旋—空间对称性。
\item 按 I3、Y\_flavor=B+S、S 组织相关矩阵。
\item 比较同多重态质量劈裂并做夸克质量外推。
\end{enumerate}\begin{columns}[T,onlytextwidth]
\column{0.56\textwidth}\begin{block}{停止条件与交叉检查}\scriptsize\begin{itemize}
\item[\CheckMark] uuu 只能位于全对称十重态。
\item[\CheckMark] uds 可同时出现在 singlet、octet 和 decuplet，需其他对称性区分。
\end{itemize}\end{block}
\column{0.40\textwidth}\begin{block}{代码映射}
\scriptsize \texttt{课程内纸笔/\allowbreak{}符号计算练习}
\end{block}\end{columns}\end{frame}

\begin{frame}[t]{夸克模型、flavor 多重态与强子分类：练习与完整解答}\FrameAnchor{23.02-5}
\footnotesize
\begin{block}{\ExerciseID{23.02}}核对三夸克 flavor 分解的总维数。\end{block}
\begin{exampleblock}{\SolutionID{23.02}}右侧维数 10+8+8+1=27，等于 3\ensuremath{^{3}}。\end{exampleblock}
\vfill
\scriptsize 回看：\CourseRef{Def}{23.02-8fa457f743}；\CourseRef{Eq}{23.02}；\CourseRef{SYM}{23.02}；\CourseRef{Alg}{23.02}。
\SourceLine{BOOK-QFT；BOOK-GROUP；BOOK-LQCD}
\end{frame}

\begin{frame}[t]{角动量耦合、宇称、荷共轭与 JPC：问题与图像}\FrameAnchor{23.03-1}
\footnotesize
\begin{block}{本节问题}给定轨道 L 和总自旋 S，哪些 JPC 能由简单 q qbar 取得？\end{block}
\PhysicalFlow{先用三角规则列出 J}{再由 L 求 P}{对中性自共轭 flavor 再求 C}{23.03}
\begin{block}{物理原则}\KnowledgeID{23.03}\quad C 只适用于相应中性、自共轭 flavor 组合；带电强子没有 C 本征值。\end{block}
\SourceLine{BOOK-QFT；PYQCD-CORRELATOR}
\end{frame}

\begin{frame}[t]{角动量耦合、宇称、荷共轭与 JPC：定义与主公式}\FrameAnchor{23.03-2}
\footnotesize\begin{tabularx}{\linewidth}{p{0.30\linewidth}Y}
\toprule 术语 & 本课程中的精确定义\\\midrule
\DefinitionID{23.03-e7dbde5a52} 角动量耦合 & 两个角动量 L、S 的张量积含 J=|L-S|,…,L+S\\
\DefinitionID{23.03-da7ebf6d83} 宇称 P & 空间反演本征值；q qbar 还含一对相反的内禀宇称\\
\DefinitionID{23.03-05de66adb5} 荷共轭 C & 仅对中性自共轭 flavor 组合定义的粒子—反粒子交换本征值\\
\bottomrule\end{tabularx}\vspace{0.15cm}
\begin{block}{\EquationID{23.03} 主公式}\centering\CourseDisplayEquation{|L-S|\le J\le L+S,\qquad P_{q\bar q}=(-1)^{L+1},\qquad C_{q\bar q}=(-1)^{L+S}}\end{block}
反夸克的内禀宇称相对夸克给额外负号；荷共轭含空间与自旋交换。
\end{frame}

\begin{frame}[t]{角动量耦合、宇称、荷共轭与 JPC：推导与证据边界}\FrameAnchor{23.03-3}
\footnotesize
\DerivationID{23.03}\quad 由本节定义与前置结论逐步推出 \CourseRef{Eq}{23.03}：
\begin{enumerate}
\item 角动量加法给 J 从 |L-S| 到 L+S，每次增加 1；q qbar 的 S 只能为 0 或 1。
\item 相对轨道波函数在 r\ensuremath{\rightarrow}-r 下给 (-1)\ensuremath{^{L}}，费米子—反费米子内禀宇称乘积为 -1，所以 P=(-1)\ensuremath{^{L+1}}。
\item 交换 q 与 qbar 时轨道给 (-1)\ensuremath{^{L}}，自旋 singlet/\allowbreak{}triplet 给 (-1)\ensuremath{^{S}}，故中性自共轭态 C=(-1)\ensuremath{^{L+S}}。
\end{enumerate}\begin{block}{可执行检查}
\SympyID{23.03}\quad 角动量耦合、宇称、荷共轭与 JPC；2/2 项通过；SymPy exact。
\par\scriptsize 假设：中性 q qbar、L 与 S 为非负整数
\end{block}
\BoundaryLine{不验证 flavor 相位和非自共轭强子的 C 标签。}
\end{frame}

\begin{frame}[t]{角动量耦合、宇称、荷共轭与 JPC：计算算法}\FrameAnchor{23.03-4}
\footnotesize
\AlgorithmID{23.03}\begin{enumerate}
\item 列出算符的空间导数阶数和 Dirac 结构。
\item 在连续群与格点 little group 下投影 J/\allowbreak{}irrep。
\item 由变换算符显式检查 P、C。
\item 用交叉相关近零和谱重叠验证不同量子数不混合。
\end{enumerate}\begin{columns}[T,onlytextwidth]
\column{0.56\textwidth}\begin{block}{停止条件与交叉检查}\scriptsize\begin{itemize}
\item[\CheckMark] L=0,S=0 给 0\ensuremath{^{-+}} 赝标介子。
\item[\CheckMark] L=0,S=1 给 1\ensuremath{^{--}} 矢量介子。
\end{itemize}\end{block}
\column{0.40\textwidth}\begin{block}{代码映射}
\scriptsize \texttt{课程内纸笔/\allowbreak{}符号计算练习}
\end{block}\end{columns}\end{frame}

\begin{frame}[t]{角动量耦合、宇称、荷共轭与 JPC：练习与完整解答}\FrameAnchor{23.03-5}
\footnotesize
\begin{block}{\ExerciseID{23.03}}简单 q qbar 能否给 JPC=1\ensuremath{^{-+}}？\end{block}
\begin{exampleblock}{\SolutionID{23.03}}P=- 要求 L 偶；C=+ 要求 L+S 偶。若 L=0 则 S=0 只能 J=0；更高偶 L 与 S\ensuremath{\le}1 也不能给 J=1 同时满足，故属 exotic。\end{exampleblock}
\vfill
\scriptsize 回看：\CourseRef{Def}{23.03-e7dbde5a52}；\CourseRef{Eq}{23.03}；\CourseRef{SYM}{23.03}；\CourseRef{Alg}{23.03}。
\SourceLine{BOOK-QFT；PYQCD-CORRELATOR}
\end{frame}

\begin{frame}[t]{两体相空间、散射截面与衰变宽度：问题与图像}\FrameAnchor{23.04-1}
\footnotesize
\begin{block}{本节问题}理论振幅怎样连同通量和相空间变成实验率？\end{block}
\PhysicalFlow{S 矩阵元 M}{Lorentz 不变相空间}{通量与对称因子}{23.04}
\begin{block}{物理原则}\KnowledgeID{23.04}\quad 相同末态粒子需除以阶乘；自旋/\allowbreak{}颜色求和平均约定必须明写。\end{block}
\SourceLine{BOOK-QFT}
\end{frame}

\begin{frame}[t]{两体相空间、散射截面与衰变宽度：定义与主公式}\FrameAnchor{23.04-2}
\footnotesize\begin{tabularx}{\linewidth}{p{0.30\linewidth}Y}
\toprule 术语 & 本课程中的精确定义\\\midrule
\DefinitionID{23.04-585297dece} 散射振幅 & 从 S 矩阵元中提出整体四动量 \ensuremath{\delta} 函数后留下的动力学量 M\\
\DefinitionID{23.04-83fe738a2c} Lorentz 不变相空间 & 每个末态含 d\ensuremath{^{3}}p/\allowbreak{}[(2\ensuremath{\pi})\ensuremath{^{3}}2E] 并由四动量 \ensuremath{\delta} 函数约束的积分\\
\DefinitionID{23.04-76c3f2092e} 通量因子 & 把单位体积入射态归一化换成单位时间碰撞率的分母\\
\bottomrule\end{tabularx}\vspace{0.15cm}
\begin{block}{\EquationID{23.04} 主公式}\centering\CourseDisplayEquation{\frac{d\sigma_{2\to2}}{d\Omega}=\frac{1}{64\pi^2s}\frac{|\bm p_f^*|}{|\bm p_i^*|}\overline{|\mathcal M|^2},\qquad d\Gamma=\frac{1}{2M}|\mathcal M|^2d\Phi_n}\end{block}
第一式是质心系两体截面；横线表示初态自旋/\allowbreak{}颜色平均与末态求和。第二式把单粒子初态换成静止质量 M 的衰变。
\end{frame}

\begin{frame}[t]{两体相空间、散射截面与衰变宽度：推导与证据边界}\FrameAnchor{23.04-3}
\footnotesize
\DerivationID{23.04}\quad 由本节定义与前置结论逐步推出 \CourseRef{Eq}{23.04}：
\begin{enumerate}
\item 相对论归一化给 d\ensuremath{\Phi}2=(2\ensuremath{\pi})\ensuremath{^{4}}\ensuremath{\delta}4(pi-pf)\ensuremath{\Pi}r d\ensuremath{^{3}}pr/\allowbreak{}[(2\ensuremath{\pi})\ensuremath{^{3}}2Er]。
\item 以质心系 \ensuremath{\delta} 函数先做总动量和末态径向积分，得到 d\ensuremath{\Phi}2/\allowbreak{}d\ensuremath{\Omega}=|pf*|/\allowbreak{}(16\ensuremath{\pi}\ensuremath{^{2}}sqrt(s))。
\item 两粒子入射通量在质心系为 4|pi*|sqrt(s)；相除得到 d\ensuremath{\sigma}/\allowbreak{}d\ensuremath{\Omega}=|pf*|overline\{|M|\ensuremath{^{2}}\}/\allowbreak{}[64\ensuremath{\pi}\ensuremath{^{2}}s|pi*|]。衰变初态没有碰撞通量，留下 1/\allowbreak{}(2M)。
\end{enumerate}\begin{block}{可执行检查}
\SympyID{23.04}\quad 两体相空间、散射截面与衰变宽度；3/3 项通过；SymPy exact。
\par\scriptsize 假设：自然单位 hbar=c=1；质心系无质量 2->2 且末态可区分；amplitude 与角度无关的教学模型
\end{block}
\BoundaryLine{验证通量/\allowbreak{}两体相空间给出的 d\ensuremath{\sigma}/\allowbreak{}d\ensuremath{\Omega} 与积分；自旋颜色平均、相同粒子因子、质量阈值和真实振幅需逐过程处理。}
\end{frame}

\begin{frame}[t]{两体相空间、散射截面与衰变宽度：计算算法}\FrameAnchor{23.04-4}
\footnotesize
\AlgorithmID{23.04}\begin{enumerate}
\item 生成所有外腿四动量并检查壳条件。
\item 计算振幅，明确初态平均和末态求和。
\item 用相空间映射做积分并包含相同粒子因子。
\item 以软/\allowbreak{}共线极限、阈值和总概率/\allowbreak{}单位检查结果。
\end{enumerate}\begin{columns}[T,onlytextwidth]
\column{0.56\textwidth}\begin{block}{停止条件与交叉检查}\scriptsize\begin{itemize}
\item[\CheckMark] 阈值以下相空间为空，宽度为零。
\item[\CheckMark] [Gamma]=能量，自然单位下寿命为能量倒数。
\end{itemize}\end{block}
\column{0.40\textwidth}\begin{block}{代码映射}
\scriptsize \texttt{课程内纸笔/\allowbreak{}符号计算练习}
\end{block}\end{columns}\end{frame}

\begin{frame}[t]{两体相空间、散射截面与衰变宽度：练习与完整解答}\FrameAnchor{23.04-5}
\footnotesize
\begin{block}{\ExerciseID{23.04}}质量可忽略的 2\ensuremath{\rightarrow}2 过程若 M 为与角度无关的常数，求总截面。\end{block}
\begin{exampleblock}{\SolutionID{23.04}}此时 |pf*|=|pi*|，d\ensuremath{\sigma}/\allowbreak{}d\ensuremath{\Omega}=|M|\ensuremath{^{2}}/\allowbreak{}(64\ensuremath{\pi}\ensuremath{^{2}}s)；积分 4\ensuremath{\pi} 得 \ensuremath{\sigma}=|M|\ensuremath{^{2}}/\allowbreak{}(16\ensuremath{\pi}s)，并须再乘任何相同粒子对称因子。\end{exampleblock}
\vfill
\scriptsize 回看：\CourseRef{Def}{23.04-585297dece}；\CourseRef{Eq}{23.04}；\CourseRef{SYM}{23.04}；\CourseRef{Alg}{23.04}。
\SourceLine{BOOK-QFT}
\end{frame}

\begin{frame}[t]{显式卷积、硬散射因子化与部分子模型：问题与图像}\FrameAnchor{23.05-1}
\footnotesize
\begin{block}{本节问题}符号 f\ensuremath{\otimes}C 究竟代表哪个变量上的什么积分？\end{block}
\PhysicalFlow{先写 Mellin 卷积的变量和边界}{再分离 hard/\allowbreak{}collinear/\allowbreak{}soft 区}{以 \ensuremath{\mu}、\ensuremath{\zeta} 演化抵消人为分隔}{23.05}
\begin{block}{物理原则}\KnowledgeID{23.05}\quad 因子化是带适用域的理论陈述；低 Q、非包容观测或 Glauber 区可能需要修正。\end{block}
\SourceLine{BOOK-QFT；P18；P31；P32}
\end{frame}

\begin{frame}[t]{显式卷积、硬散射因子化与部分子模型：定义与主公式}\FrameAnchor{23.05-2}
\footnotesize\begin{tabularx}{\linewidth}{p{0.30\linewidth}Y}
\toprule 术语 & 本课程中的精确定义\\\midrule
\DefinitionID{23.05-5d83731ffd} Mellin 卷积 & (f\ensuremath{\otimes}C)(x)=\ensuremath{\int}x\ensuremath{^{1}} dy f(y)C(x/\allowbreak{}y)/\allowbreak{}y\\
\DefinitionID{23.05-27c38feaa3} 因子化 & 把主导动量区域组织为通用长程函数与过程依赖短程核的带适用域定理\\
\DefinitionID{23.05-303fe7623c} 尺度对 & \ensuremath{\mu} 控制普通 UV/\allowbreak{}collinear 演化；TMD 另需 rapidity 尺度 \ensuremath{\zeta}，并与 soft subtraction 方案共同定义\\
\bottomrule\end{tabularx}\vspace{0.15cm}
\begin{block}{\EquationID{23.05} 主公式}\centering\CourseDisplayEquation{(f\otimes C)(x)=\int_x^1\frac{dy}{y}f(y,\mu)C\!\left(\frac{x}{y},Q,\mu\right),\qquad \sigma=\sum_{ij}f_i\otimes f_j\otimes\hat\sigma_{ij}+O((\Lambda/Q)^p)}\end{block}
\ensuremath{\mu} 依赖在 PDF 演化与短程核之间相消到所用微扰阶；剩余依赖估计截断误差。
\end{frame}

\begin{frame}[t]{显式卷积、硬散射因子化与部分子模型：推导与证据边界}\FrameAnchor{23.05-3}
\footnotesize
\DerivationID{23.05}\quad 由本节定义与前置结论逐步推出 \CourseRef{Eq}{23.05}：
\begin{enumerate}
\item 母部分子携带动量分数 y，硬核使用子分数 z=x/\allowbreak{}y；Jacobian 给 dy/\allowbreak{}y，且 0<x\ensuremath{\le}y\ensuremath{\le}1。
\item 按 loop 动量区域分离 hard、collinear 和 soft 主导贡献；对 hard 子图短距展开，把长程敏感性吸收进规范不变分布。
\item 用明确的 zero-bin/\allowbreak{}soft subtraction 去除区域重叠；RG 与 Collins--Soper 方程分别控制 \ensuremath{\mu}、\ensuremath{\zeta}，使物理截面对人为尺度的依赖在所用阶相消。
\end{enumerate}\begin{block}{可执行检查}
\SympyID{23.05}\quad 显式卷积、硬散射因子化与部分子模型；2/2 项通过；SymPy exact。
\par\scriptsize 假设：p>0 且固定 Lambda；0<x<1；教学卷积取 f(y)=y、C(x/\allowbreak{}y)=x/\allowbreak{}y
\end{block}
\BoundaryLine{只验证一个显式卷积和高尺度幂修正；QCD 因子化的区域分离、Glauber 条件与核的微扰计算不是 CAS 命题。}
\end{frame}

\begin{frame}[t]{显式卷积、硬散射因子化与部分子模型：计算算法}\FrameAnchor{23.05-4}
\footnotesize
\AlgorithmID{23.05}\begin{enumerate}
\item 先明确过程、观测量和所需 collinear/\allowbreak{}TMD 因子化定理。
\item 选共同 \ensuremath{\mu}、\ensuremath{\zeta} 并准备所有 flavor/\allowbreak{}胶子分布。
\item 计算同阶硬核、演化和卷积。
\item 扫描尺度、核阶、幂修正和和规则，并与数据相关拟合。
\end{enumerate}\begin{columns}[T,onlytextwidth]
\column{0.56\textwidth}\begin{block}{停止条件与交叉检查}\scriptsize\begin{itemize}
\item[\CheckMark] 若 C(z)=\ensuremath{\delta}(1-z)，卷积精确退化为 f(x)。
\item[\CheckMark] Q\ensuremath{\rightarrow}\ensuremath{\infty} 时幂修正趋零，但大对数仍需 RG/\allowbreak{}CS 演化。
\end{itemize}\end{block}
\column{0.40\textwidth}\begin{block}{代码映射}
\scriptsize \texttt{课程内纸笔/\allowbreak{}符号计算练习}
\end{block}\end{columns}\end{frame}

\begin{frame}[t]{显式卷积、硬散射因子化与部分子模型：练习与完整解答}\FrameAnchor{23.05-5}
\footnotesize
\begin{block}{\ExerciseID{23.05}}证明树级核 C(z)=\ensuremath{\delta}(1-z) 时 (f\ensuremath{\otimes}C)(x)=f(x)。\end{block}
\begin{exampleblock}{\SolutionID{23.05}}根 1-x/\allowbreak{}y=0 在 y=x，且 \ensuremath{\delta}(1-x/\allowbreak{}y)=x\ensuremath{\delta}(y-x)；代入 \ensuremath{\int}x\ensuremath{^{1}}dy f(y)x\ensuremath{\delta}(y-x)/\allowbreak{}y=f(x)（端点按同一分布约定处理）。\end{exampleblock}
\vfill
\scriptsize 回看：\CourseRef{Def}{23.05-5d83731ffd}；\CourseRef{Eq}{23.05}；\CourseRef{SYM}{23.05}；\CourseRef{Alg}{23.05}。
\SourceLine{BOOK-QFT；P18；P31；P32}
\end{frame}

\begin{frame}[t]{第 23 卷能力清单}\FrameAnchor{V23-outcomes}
\small\begin{itemize}
\item[\CheckMark] 能读标准模型量子数表
\item[\CheckMark] 能用 JPC 和 flavor 分类强子
\item[\CheckMark] 能连接散射振幅、截面和部分子因子化
\end{itemize}\vfill\begin{block}{通过标准}
不以“看懂”为标准：应能从空白纸推导主公式、实现算法、解释检查失败意味着什么，并完成本卷练习。
\end{block}\end{frame}

\begin{frame}[t]{第 23 卷来源（1/2）}\FrameAnchor{V23-sources}
\scriptsize\begin{tabularx}{\linewidth}{p{0.17\linewidth}p{0.28\linewidth}Y}
\toprule ID & 来源 & 本卷用途\\\midrule
\SourceID{BOOK-QFT} & An Introduction to Quantum Field Theory（仓库转排本） & 量子场论、规范理论、QCD 和重整化的主教材交叉核验。\\
\SourceID{BOOK-GROUP} & 连续群与生成元预备课（课程自编） & 从旋转群过渡到 SU(2)/\allowbreak{}SU(3)。\\
\SourceID{BOOK-LQCD} & Introduction to Lattice QCD /\allowbreak{} 格点 QCD 导论（仓库转排本） & 欧氏化、规范链接、费米子、连续极限和关联函数。\\
\SourceID{PYQCD-CORRELATOR} & PyQCD 关联函数技能 & 强子二点/\allowbreak{}三点与谱学检查。\\
\SourceID{P18} & 大动量有效理论综述 & 论文图谱 P18 的完整原始 PDF；底稿 arXiv:2004.03543；SHA-256 bfcdc7b4100d0da2…。\\
\bottomrule\end{tabularx}\end{frame}

\begin{frame}[t]{第 23 卷来源（2/2）}\FrameAnchor{V23-sources-2}
\scriptsize\begin{tabularx}{\linewidth}{p{0.17\linewidth}p{0.28\linewidth}Y}
\toprule ID & 来源 & 本卷用途\\\midrule
\SourceID{P31} & CT18全局分析 & 论文图谱 P31 的完整原始 PDF；底稿 arXiv:1912.10053；SHA-256 98d4b253c3c926d5…。\\
\SourceID{P32} & NNPDF17高精度数据部分子分布 & 论文图谱 P32 的完整原始 PDF；底稿 arXiv:1706.00428；SHA-256 9d3daadd61fc48bc…。\\
\bottomrule\end{tabularx}\end{frame}

